% ================= APPENDIX: RISK REGISTER AND CONTINGENCY PLAN =================
\section{Risk Register and Contingency Plan}
\label{app:risks}

Table~\ref{tab:risk_register} lists the risks judged material to the project,
with the trigger that realises each one, its impact, severity (S) and
likelihood (L) on an H/M/L scale, the accountable owner, and the mitigation.
Table~\ref{tab:cascade} fixes, in advance, the order scope is cut if the
schedule runs short, and the note that follows covers how the team backs up
whoever is stuck on a critical-path task.

\begin{footnotesize}
\setlength{\tabcolsep}{4pt}
\begin{longtable}{@{}L{0.8cm} L{2.0cm} L{2.3cm} L{2.3cm} c c L{1.2cm} L{3.5cm}@{}}
\caption[Risk register]{Risk register: trigger, impact, severity (S),
likelihood (L), owner and mitigation.}
\label{tab:risk_register}\\
\toprule
\textbf{ID} & \textbf{Risk} & \textbf{Trigger} & \textbf{Impact} & \textbf{S} & \textbf{L} & \textbf{Owner} & \textbf{Mitigation} \\
\midrule\endfirsthead
\toprule
\textbf{ID} & \textbf{Risk} & \textbf{Trigger} & \textbf{Impact} & \textbf{S} & \textbf{L} & \textbf{Owner} & \textbf{Mitigation} \\
\midrule\endhead
\bottomrule\endfoot
\multicolumn{8}{@{}l}{\textbf{Electrical and Power}} \\
\midrule
RE1 & \qty{16}{\volt} capacitor on the \qty{25}{\volt} bus & Build error: the BOM \qty{100}{\micro\farad} parts are \qty{16}{\volt} and fit the bus rail & Capacitor vents; board destroyed and bus shorted; ${\ge}\,3$~day recovery through the 9~Aug gate & H & M & Pablo & \qty{16}{\volt} parts segregated at build, \qty{5}{\volt} rail only; bulk cap is a separate \qtyrange{470}{1000}{\micro\farad}, ${\ge}\,\qty{50}{\volt}$ item; checked before any power \\
RE2 & No bulk capacitance at regenerative braking & Hard wheel deceleration returns ${\approx}\,\qty{72.8}{\joule}$ per wheel (\qty{218}{\joule} across three) to the bus & Bus overvoltage damages a moteus-n1; loses one axis and the bench-rig spare & H & M & Pablo & Fit and verify the \qty{50}{\volt} bulk capacitor before any commanded deceleration; driver bus-overvoltage fault limit \\
RE3 & MA600 air gap out of specification & Gap ${>}\,\qty{0.5}{\milli\metre}$ or chip off-centre, at build or disturbed later at assembly & Angle dropout and degraded commutation; presents as a control fault and misdirects debugging & H & H & Pablo & Shim value fixed on the bench and re-verified per axis at every assembly step; encoder error flags logged in every balance run \\
RE4 & SPI run exceeds \qty{20}{\centi\metre} & Repackaging moves a driver away from its encoder & Encoder corruption at \qty{6}{\mega\hertz}; same signature as RE3 & H & M & Pablo & Driver positions derived from encoder positions and frozen in the board outline; later moves are electrical change requests \\
RE5 & CAN-FD unreliable at \qty{5}{\mega\bit\per\second} & Stub ${>}\,\qty{10}{\centi\metre}$, mis-placed termination, or star topology at integration & Error frames and dropped commands under vibration; intermittent and hard to diagnose & H & M & Pablo & Linear topology; split termination at the two physical bus ends only; scope and soak test before integration \\
RE6 & Board outline late to CAD & Layout slips past 1~Aug & Mass freeze and frame print slip, cascading to the 9~Aug integration gate & H & M & Pablo & Two days of float; if bring-up slips, issue the outline with conservative keep-outs rather than delay it \\
RE7 & Rig parameters late to control & Wheel balance or rig assembly slips past 2~Aug & Blocks the sequential control chain; propagates to 20~Aug with no recovery & H & M & Pablo & Escalate if off track; fallback delivers $K_t$ and loop latency from bench measurement without the wheel \\
RE8 & \qty{5}{\volt} rail undersized & Teensy, XIAO, IMU and servo stall draw together & Rail sag and MCU brown-out mid-balance; misread as a control failure & M & M & Pablo & Sum the load budget before build; heatsink the regulator; bulk cap plus back-feed diode; measure sag under worst-case load \\
RE9 & Battery bring-up incident & Battery connected before the bench sequence completes, or without anti-spark switch and fuse & Arc or LiPo fire; injury and total project loss & H & L & Pablo & All bring-up on a current-limited supply; battery introduced only at the final power-on, with anti-spark switch, fuse, cell alarm and a second person present \\
RE10 & Power harness fails under load & Cold joint or under-gauge wire on a high-current bus & Voltage drop, connector melt, or open circuit mid-demonstration & M & M & Pablo & \qty{14}{\awg} minimum, pull-tested and thermally imaged; spare harness in the demonstration kit \\
RE11 & No spare actuation set & Any motor, driver or encoder fails after cube integration & Loses an axis with no resupply lead time before 20~Aug; demonstration cut short & H & M & Pablo & All three sets treated as flight hardware from integration on; no cannibalising for bench tests afterward; reseat and inspect at every touch \\
RE12 & Brake work displaces edge balance & Brake build runs concurrently with the edge-balance gate & Missing the gate cuts the jump-up stretch goal, which the brake exists to serve & M & M & Pablo & Brake support is the designated sacrificial task; conflicts resolve to edge balance by standing decision \\
RE13 & Vibration loosens connectors & Sustained running at high wheel speed & Intermittent axis dropout; cube falls and damages structure & M & H & Pablo & Strain relief, labelling and connector retention; reseat check as the first action of every support call \\
RE14 & Perfboard rework consumes its float & Wiring error found at bus test; the build step has no slack & Board late to integration, cascading to the 9~Aug gate & M & M & Pablo & Wiring checked against the floorplan before power; rails brought up incrementally so faults localise \\
\multicolumn{8}{@{}l}{\textbf{Programme and Technical}} \\
\midrule
PR1 & Final week has almost no slack & Any slip after 8~Aug & No days left to recover before the 20~Aug presentation & H & M & Pablo & Go/no-go on the jump-up called at end of day 11~Aug; nothing new started after that date, only debugging what already runs \\
PR2 & Structure heavier than the mass allowance & Printed frame and fasteners weigh in over the \qty{374}{\gram} allowance & Wheel undersized against the jump-up sizing; momentum budget missed & M & M & Nasia & \qty{181}{\gram} of headroom before re-ballasting is needed; weigh at frame print and re-run the sizing script if over \\
PR3 & Brake torque never measured on the bench & Brake work stays sacrificial through the final week & The inertia-requirement safety margin ships unverified & M & H & Niccolo & Run a standalone spin-down test in any open window; otherwise state explicitly that the figure is unmeasured \\
PR4 & Programme lead is a single point of failure & Pablo unavailable during a critical-path week & Schedule, most electrical tasks and several software tasks stall together & H & L & All & Deyan holds the document and schedule state; Niccolo and Andrea can act on the electrical notes directly; no single day depends on one person \\
\end{longtable}
\end{footnotesize}

\begin{table}[htbp]
\centering
\caption[Contingency: order of cuts]{Contingency: order in which scope is cut
if the schedule runs short.}
\label{tab:cascade}
\small
\begin{tabular}{@{}c L{9.5cm}@{}}
\toprule
\textbf{Order} & \textbf{Scope cut} \\
\midrule
1 & Jump-up demonstration (S3) --- cut first; standing decision, see \S\ref{sec:schedule} \\
2 & Brake mechanism build and its bench test --- cut with the jump-up, since it exists only to serve it \\
3 & Corner balance and commanded slew (S2b/c) --- cut only if edge balance is not solid by 17~Aug \\
4 & Edge balance (S2a) and the 1-DoF rig result --- never cut; the protected demonstration and the document's primary evidence \\
\bottomrule
\end{tabular}
\end{table}

\paragraph{Parallel work and backup coverage.} The five workstreams
(documentation, control/software, hardware build, 3D design, electrical) run
in parallel by design, so a stall in one does not by default stall the
others --- but only if someone is deliberately free to help. The practice
already applies twice in this plan: Nasia's fallback on any day the sizing
numbers are late is CAD review and tolerance work rather than idle waiting,
and Andrea's fallback if the rig is not ready on time is bench measurement of
$K_t$ and loop latency, so she is never simply blocked (RE7). The same rule
holds for the rest of the schedule: whoever is not the bottleneck on a given
day supports whoever is, rather than working ahead alone on their own next
task.
